\chapter{Ștergerea din structuri de date care nu admit ștergeri}

Există situații cînd avem nevoie de o structură de date $S$ care să suporte adăugări și ștergeri. Știm să facem adăugarea într-o complexitate $\bigoh(T(n))$, dar nu știm deloc să facem ștergerea. Să spunem, în schimb, că știm să facem operații \textit{undo} (ștergerea elementelor în ordinea inversă adăugării lor). Această operație este adesea mai simplă. În această situație, putem folosi arborii de intervale ca să implementăm și ștergeri în general, dacă plătim un cost suplimentar de $\bigoh(\log n)$. Așadar, complexitatea finală va fi $\bigoh(T(n) \log n)$ per operație.

Iată un exemplu. Să spunem că dorim o structură care să admită adăugări, ștergeri și interogări de maxim. Am putea folosi un \ccode{set} din STL (pe care noi înșine nu-l putem implementa ușor). Dar, dacă mereu am avea de șters doar ultimul element adăugat, am putea implementa structura absolut elementar! Cum?

\begin{turn}{180}
  \begin{minipage}{\linewidth}
    Folosim o singură stivă în care ținem perechi de valori: ultima valoare adăugată și maximul pînă la acel moment. Pentru a șterge un element, eliminăm un element din stivă, iar maximul din noul vîrf al stivei este maximul elementelor rămase.
  \end{minipage}
\end{turn}

În supa culturală românească, tehnica se numește „AINT pe timp”. \href{https://cp-algorithms.com/data_structures/deleting_in_log_n.html}{CP Algorithms} are un articol foarte bun despre subiect.
